\section{本章小结}

本章我们完整拆解了 GDT 的每一个字节：

\begin{enumerate}
  \item \textbf{段选择子}（16 位）= index + TI + RPL
  \item \textbf{描述符}（8 字节）= 碎片化的 base/limit + access byte + flags
  \item \textbf{段寄存器}的双面结构：可见的 selector + 不可见的 hidden cache
  \item \textbf{gdt\_flush.asm}：lgdt 加载表 → far jump 刷新 CS → mov 刷新 DS/SS
\end{enumerate}

GDT 是 x86 保护模式的基础设施，也是后续 TSS（Ring 3 栈切换）和
IDT（中断门需要代码段 selector）的前置条件。

下一章我们设置 IDT，让内核能捕获 CPU 异常并打印诊断信息。

